\begin{description}
\item[(a)] Work out the mass density of the PBH,
  \[ \rho_{\mathrm{PBH}} = \frac{M}{V_{\mathrm{S}}}
     = \frac{M}{\frac{4}{3}\pi r_{\mathrm{S}}^3} \]
  \hfill(1 point)

  where the Schwarzschild radius is
  \[ r_{\mathrm{S}} = \frac{2GM}{c^2} \]
  \hfill(0.5 points)

  The above equation becomes
  \[ \rho_{\mathrm{PBH}}
     = \frac{M}{\frac{4}{3}\pi\left(\frac{2GM}{c^2}\right)^3}
     = \frac{3c^6}{32\pi G^3 M^2} \]
  \hfill(0.5 points)

  The PBHs can be formed if the cosmological density of the universe is of
  the same order as the density of the PBH,
  \[ \rho_{\mathrm{c}} \sim \rho_{\mathrm{PBH}} \]
  \hfill(2 points)

  Then, we can figure out the desired relation,
  \[ \frac{1}{Gt^2} \sim \frac{3c^6}{32\pi G^3 M^2}
     \Rightarrow M \sim \frac{c^3 t}{6G} \]
  \hfill(1 point)

  \textit{Note:}
  \begin{itemize}
  \item \textit{If students use ``$=$'', ``$\approx$'', or ``$\le$'' for the
    density relation above, the marks remain unchanged.}
  \item \textit{Students can write the coefficient as $\sqrt{3/32\pi}$ or 1
    instead of 1/6.}
  \end{itemize}

\item[(b)] The PBHs still evaporating implies that the evaporation time must
  be greater than or equal to the Universe age,
  \[ \tau(M) \ge 13.8\ \mathrm{(Gyr)} \]
  \hfill(1.5 points)
  \begin{align*}
    &\Rightarrow 10^{67}\left(\frac{M}{M_\odot}\right)^3 \ge 13.8 \times 10^9 \\
    &\Rightarrow \left(\frac{M}{M_\odot}\right)^3 \ge 13.8 \times 10^{-58} \\
    &\Rightarrow M \ge 1.1 \times 10^{-19} \times M_\odot
  \end{align*}
  \hfill(1 point)

  Then, using the result derived in part a), we can work out the formation
  time after the Big Bang of PBHs as follows
  \[ \frac{c^3 t}{6G} \gtrsim 1.1 \times 10^{-19} \times M_\odot \]
  \hfill(1.5 points)
  \begin{align*}
    &\Rightarrow t \gtrsim 6.6 \times 10^{-19} \times M_\odot \times \frac{G}{c^3} \\
    &\Rightarrow t \gtrsim 3.3 \times 10^{-24}\ \mathrm{(s)}
  \end{align*}
  \hfill(1 point)

  \textit{Note:}
  \begin{itemize}
  \item \textit{Deduct 0.5 points if students use either
    ``$\tau(M) = 13.8\,\mathrm{Gyr}$'' or ``$\tau(M) > 13.8\,\mathrm{Gyr}$''.}
  \item \textit{The above comparisons are in the same order of magnitude.
    Therefore, coefficients are not important or can be flexibly varied.}
  \end{itemize}
\end{description}
